\section{Finite relational presentations and witness amalgamation}
\label{sec:relations}

\subsection{Layered candidates}

A finite relational presentation permits nonunique answers.  Fix a parameter
\(p\).  At level \(k\), candidates are numbers \(a<B(p,k)\).  Primitive-
recursive bits assert
\[
 \Good(p,k,a),
 \qquad
 \Compat(p,k,a,b),
\]
where the second relates level \(k\) to level \(k+1\).  A finite path is a
coded list whose entries are in range and good and whose adjacent pairs are
compatible.

The central issue is dead ends.  Choosing the least good candidate at each
level can fail even when a complete path exists.  Backward viability removes
exactly the candidates that cannot reach the requested horizon.

\subsection{Backward viability}

Define viability at remaining depth \(d\) by
\begin{align*}
 \Viable(p,k,0,a)&=[a<B(p,k)]\wedge\Good(p,k,a),\\
 \Viable(p,k,d+1,a)&=[a<B(p,k)]\wedge\Good(p,k,a)\wedge{}\\
 &\quad\Any_{b<B(p,k+1)}
 \bigl(\Good(p,k+1,b)\wedge\Compat(p,k,a,b)
       \wedge\Viable(p,k+1,d,b)\bigr).
\end{align*}
All operations are Boolean folds, so this is a fixed primitive-recursive
predicate.

\begin{theorem}[Viability/path equivalence]
\label{thm:viable-path}
iPRA proves that \(\Viable(p,k,d,a)=1\) exactly when there exists a coded good
compatible path of \(d\) edges beginning with \(a\) at level \(k\).  The path
code has a primitive-recursive bound determined by \(B,p,k,d\).
\end{theorem}

\begin{proof}
Induct on \(d\).  At zero, the singleton path condition is the definition.  At
a successor, fold correctness extracts a compatible viable successor \(b\).
Apply the induction hypothesis to the tail and prepend \(a\).  Conversely,
split a path into its first edge and tail and use the same bounded witness in
the viability fold.  For the formal induction, compare the viability bit with
a bounded path-existence checker; the equality of bits is quantifier free.
\end{proof}

\subsection{Canonical finite paths}

When \(a\) is viable at depth \(d+1\), bounded least search returns the least
successor \(b\) that is good, compatible, and viable at depth \(d\).  Repeating
this search by primitive recursion yields a canonical path.

\begin{corollary}[Least viable path]
\label{cor:least-path}
There is a fixed primitive-recursive path constructor \(\mathsf{path}\) such
that iPRA proves
\[
 \Viable(p,k,d,a)=1\to
 \mathsf{PathOK}(p,k,d,a,\mathsf{path}(p,k,d,a))=1.
\]
\end{corollary}

\begin{proof}
Induct on the generated prefix length.  The invariant states that the current
node is viable for the remaining depth and that the accumulated prefix passes
the path checker.  Least-search correctness supplies the next viable edge.  At
the last level, depth-zero viability gives goodness of the endpoint.
\end{proof}

This is finite dependent choice implemented by bounded search.  It has no
infinite branch and no use of countable choice.

\subsection{Full-horizon viability}

For cross-level comparison, forward extendibility alone is insufficient.  A
high-level candidate may extend indefinitely but be disconnected from the
input.  Define
\(\FullV(p,h,k,a)=1\) to mean that \(a\) occurs at level \(k\) on a complete
good path from level zero to horizon \(h\).

Let \(\mathsf{Reach}(p,k,a)\) be forward reachability from a good level-zero
candidate.  Then, for \(k\le h\),
\[
 \FullV(p,h,k,a)
 \quad\leftrightarrow\quad
 \mathsf{Reach}(p,k,a)\wedge\Viable(p,k,h-k,a).
\]
The proof splits or concatenates finite paths.  This factorization supports a
dynamic-programming implementation with a forward table and a backward table.
A direct bounded path-code checker is simpler for certificates; both
implementations have the same object specification.

\subsection{Amalgamating finite requirements}

Suppose a theorem has finitely many requirements: approximation bounds, chart
membership, algebraic identities, compatibility with a second module, and a
final residual test.  Conjoin them into the goodness and edge predicates.

\begin{theorem}[Finite witness amalgamation]
\label{thm:amalgamation}
If
\[
 \Any_{a<B(p,0)}\Viable(p,0,h,a)=1,
\]
then bounded search followed by the least viable path constructor returns one
finite witness family satisfying every compiled requirement simultaneously.
Adding a requirement can only remove viable nodes; it cannot create a spurious
path.
\end{theorem}

\begin{proof}
The initial \(\Any\)-fold and least-search theorem produce a viable root.
Apply \cref{cor:least-path}.  Soundness of Boolean conjunction exposes each
component requirement.  For monotonicity, induct on remaining depth: every
witness edge valid under the stronger goodness predicate is also valid under
the weaker one.
\end{proof}

\subsection{When a relation is the final theorem}

A finite relation is not merely an intermediate inconvenience.  It is often
the correct stable output:
\begin{itemize}
\item a root multiset or a family of isolating clusters at collisions;
\item a finite cover or parity certificate in topology;
\item approximate rotation data at the identity;
\item a collection of descent representatives with local compatibility;
\item a finite interval of possible minimizers while the minimum value is
      already stable.
\end{itemize}
Downstream statements should quantify over all accepted witnesses or consume a
stable observable.  A selector is introduced only if a theorem proves that its
choice is harmless.
